\section*{Chapitre \ref{ch:iv2:credit-and-xva} --- Crédit et XVA}
\markboth{\protect\omTocMark{}{Réponses~: Crédit et XVA}}{\protect\omTocMark{}{Réponses~: Crédit et XVA}}

\begin{solution}{iq:iv2:credit-and-xva:1}
L'intensité de défaut vaut environ 3.33~\% par an, et la probabilité de défaut sur cinq ans est d'environ 15.4~\%.

Il s'agit du \index{triangle du crédit}triangle du crédit~: le spread de crédit vaut l'intensité de défaut $\times$ (1 $-$ taux de recouvrement), donc $h = 0.02/0.6$. La probabilité de survie vaut $e^{-5h} = e^{-1/6} = 0.846$.

La logique est celle d'un échange équitable. Chaque année, la prime versée est le spread de crédit. Le versement attendu est la probabilité de défaut multipliée par la perte en cas de défaut (LGD).

\iqwrong{Multiplier $5 \times 3.33~\% = 16.7~\%$ et y voir la probabilité. C'est proche ici, mais cela ignore la survie et peut dépasser 100~\% pour un émetteur en difficulté.}
\end{solution}

\begin{solution}{iq:iv2:credit-and-xva:2}
Elle vaut 3.92~\% la première année, 3.77~\% la deuxième année, vue d'aujourd'hui, et 3.92~\% à nouveau sachant la survie à la première année.

$1 - e^{-0.04} = 0.0392$. Depuis aujourd'hui, la deuxième année suppose la survie puis le défaut~: $e^{-0.04} - e^{-0.08} = 0.0377$. Une intensité constante n'a pas de mémoire, donc le chiffre conditionnel se répète.

\iqwrong{Dire que la deuxième année doit être plus risquée parce que le temps a passé. Avec une intensité constante, ce n'est pas le cas. Le chiffre inconditionnel n'est inférieur que parce que l'entreprise a pu déjà disparaître.}
\end{solution}

\begin{solution}{iq:iv2:credit-and-xva:3}
Sans recouvrement, le contrat vaut 0.7788, soit un spread de crédit exactement égal à 3~\%. Avec un recouvrement de 40~\%, il vaut 0.8292, soit un spread de crédit d'environ 1.75~\%.

Sans recouvrement~: $e^{-(0.02 + 0.03) \times 5} = e^{-0.25}$. Le spread de crédit est l'intensité de défaut.

Avec un recouvrement égal au nominal à l'échéance~: $e^{-0.1}\,(e^{-0.15} + 0.4(1 - e^{-0.15})) = 0.9048 \times 0.9164$. Le spread de crédit vaut $-\ln(0.9164)/5 = 1.746~\%$, un peu en dessous de $0.6 \times 3~\% = 1.8~\%$.

\iqwrong{Prendre 1.8~\% pour une valeur exacte. Le triangle du crédit est le résultat valable pour une prime périodique. Pour un zéro-coupon à cinq ans avec recouvrement à l'échéance, ce n'est qu'une approximation, et l'écart se creuse avec une intensité plus élevée et une durée plus longue.}
\end{solution}

\begin{solution}{iq:iv2:credit-and-xva:4}
On paie 50\,000~\$ par trimestre, et en cas de défaut on reçoit 7~millions de dollars. On est short sur le crédit~: on gagne lorsque l'émetteur se dégrade.

La prime annuelle se calcule ainsi~: $0.02 \times 10$ millions de dollars $= \$200\,000$. Le versement est de $(1 - 0.3) \times 10$ millions de dollars. La prime courue jusqu'à la date de défaut est également réglée.

\iqwrong{Appeler un acheteur de protection «~long~», comme s'il était détenteur d'une option. Sur un desk de crédit, être long du crédit signifie porter le risque, comme détenir l'obligation. Le vendeur de protection est donc long.}
\end{solution}

\begin{solution}{iq:iv2:credit-and-xva:5}
L'acheteur de protection verse au vendeur environ 4.3 points de paiement initial. Cela fait environ 430\,000~\$ sur 10~millions de dollars.

Depuis 2009, les contrats nord-américains se négocient avec des coupons fixes de 100 ou 500 points de base, et la différence est réglée d'avance. L'acheteur paie 100 par an pour une protection valant 200, si bien que l'écart est payé d'emblée.

L'intensité vaut $h = 0.02/0.6 = 3.33~\%$. L'annuité risquée vaut $(1 - e^{-(r + h)T})/(r + h)$ avec $r + h = 6.33~\%$, soit 4.29. Le paiement initial vaut donc $(0.02 - 0.01) \times 4.29 = 4.29~\%$.

\iqwrong{Multiplier par 5 ans pour obtenir 5 points. La prime supplémentaire s'arrête au défaut et est actualisée, si bien que le facteur est l'annuité risquée, et non la durée de vie du contrat.}
\end{solution}

\begin{solution}{iq:iv2:credit-and-xva:6}
À 200 points de base, le paiement initial vaut 4.29 points avec un recouvrement de 40~\% et 4.37 avec 20~\%. À 2\,000 points de base, il vaut 43.8 points et 51.1. Le recouvrement fait donc varier le premier cas de 0.1 point et le second de plus de 7 points.

La formule est la même~: paiement initial = (spread $-$ coupon) $\times$ annuité risquée, avec $h = s/(1 - R)$.

Pour un émetteur à 200 points de base, l'intensité passe de 3.33~\% à 2.5~\%, et l'annuité ne passe que de 4.29 à 4.37. L'écart entre les spreads est faible, si bien que cela compte à peine.

Pour un émetteur à 2\,000 points de base, l'intensité passe de 33.3~\% à 25~\%. L'annuité passe alors de 2.30 à 2.69, et l'écart à multiplier vaut 19 points.

\iqwrong{Dire que le recouvrement n'a pas d'importance parce que la cotation est déjà exprimée en spread. C'est vrai pour un contrat au pair, où le paiement initial est nul quel que soit le recouvrement. Cela ne tient plus dès que le spread s'éloigne du coupon.}
\iqfollow{Comment un émetteur en difficulté serait-il coté~?}{Souvent en points de paiement initial, et parfois en prix, ce qui laisse moins matière à discussion qu'un spread qui dépend d'une hypothèse de recouvrement.}
\end{solution}

\begin{solution}{iq:iv2:credit-and-xva:7}
Acheter l'obligation et acheter une protection. Il s'agit d'un trade de base négative, qui rapporte 40 points de base par an, soit 40\,000~\$ sur 10~millions de dollars, tant qu'il dure.

\textbf{Pourquoi les spreads doivent coïncider.} On détient un floater au pair de spread $s_b$ et une protection à $s_c$, le tout étant financé au taux de financement. En l'absence de défaut, on encaisse $s_b - s_c$. En cas de défaut, la protection verse le pair moins le recouvrement, on cède l'obligation au recouvrement, et l'on rembourse le pair. On ne porte donc aucun risque de crédit, et le gain courant doit être nul si le taux de financement est le vrai coût de financement~: $s_b = s_c$. Avec une intensité $h$ et un recouvrement égal au pair, les deux valent $h(1 - R)$.

\textbf{Ce qui peut mal tourner.}
\begin{itemize}
\item Le financement devient plus cher, ou le repo est retiré. Le trade ne porte aucun risque de crédit, mais il porte un risque de financement, et les basis trades sont réduits lorsque le financement se resserre.
\item L'obligation n'est pas au pair au défaut, si bien que le versement ne correspond pas. Une obligation cotée au-dessus du pair perd davantage que ce que verse la protection.
\item Le contrat et l'obligation peuvent diverger sur ce qui constitue un défaut ou sur ce qui peut être livré.
\item Le trade peut afficher une perte avant de converger, et l'on peut être contraint de sortir à ce moment-là.
\end{itemize}

\iqwrong{Qualifier cela d'argent gratuit. Les 40 points de base supplémentaires sont la rémunération que le marché verse pour immobiliser du bilan et porter un risque de financement.}
\end{solution}

\begin{solution}{iq:iv2:credit-and-xva:8}
Les deux font défaut avec une probabilité de 1.675~\%, et au moins l'un des deux avec une probabilité de 8.325~\%. Avec des probabilités de 5~\% et 1~\%, la corrélation ne peut pas dépasser environ 0.438.

Pour les variables indicatrices de défaut, la covariance vaut $P(\text{les deux}) - pq$, et la corrélation divise celle-ci par $\sqrt{p(1 - p)q(1 - q)}$. Ainsi
\[
P(\text{les deux}) = 0.05^2 + 0.3 \times 0.05 \times 0.95 = 0.01675,
\]
et la probabilité qu'au moins l'un des deux fasse défaut vaut $0.10 - 0.01675 = 0.08325$.

Avec des probabilités inégales, la probabilité conjointe de défaut ne peut pas dépasser la plus petite, soit 1~\%. La corrélation maximale possible vaut donc
\[
\frac{0.01 - 0.0005}{\sqrt{0.05 \times 0.95 \times 0.01 \times 0.99}} = 0.438.
\]

\iqwrong{Supposer qu'une corrélation de 1 est toujours possible. Elle exige que les deux événements soient un seul et même événement, ce qui est impossible avec des probabilités différentes. Cela signifie aussi que les chiffres de «~corrélation de défaut~» ne sont pas comparables entre des émetteurs de qualité très différente.}
\iqfollow{Quelle est la corrélation minimale possible ici~?}{Sans chevauchement, la probabilité conjointe est nulle, donc la corrélation vaut $-0.0005/0.0217 = -0.023$. Des événements rares ne peuvent pas non plus être fortement corrélés négativement.}
\end{solution}

\begin{solution}{iq:iv2:credit-and-xva:9}
Cas indépendants~: environ 2.94~\% par an, proche de la somme des cinq spreads de 0.6~\% chacun. Cas parfaitement corrélés~: 0.6~\%, soit comme pour un seul émetteur. Le second-to-default évolue dans l'autre sens et augmente avec la corrélation.

La probabilité qu'il y ait au moins un défaut est $1 - 0.99^5 = 4.90~\%$, et avec un recouvrement de 40~\% la perte attendue vaut $0.6 \times 4.90~\% = 2.94~\%$. Avec une corrélation parfaite, ils font tous défaut ensemble avec une probabilité de 1~\%, si bien que la perte vaut $0.6~\%$.

Le second-to-default exige deux défauts. Dans le cas indépendant, cela donne $1 - 0.99^5 - 5 \times 0.01 \times 0.99^4 = 0.098~\%$, ce qui est rare. Lorsque la corrélation augmente, les défauts arrivent par grappes. Sous une copule gaussienne, à une corrélation de 0.6, la probabilité d'un premier défaut passe de 4.90~\% à environ 3.7~\%, soit un spread d'environ 2.2~\%, et celle d'un second défaut passe de 0.098~\% à environ 0.9~\%, soit un spread d'environ 0.55~\%.

\iqwrong{Croire que les deux trades de panier progressent avec la corrélation. Le détenteur d'un first-to-default veut que les défauts soient dispersés, car cela rend au moins un défaut plus probable. Le détenteur d'un second-to-default veut au contraire qu'ils surviennent ensemble.}
\end{solution}

\begin{solution}{iq:iv2:credit-and-xva:10}
$\mathrm{EE}(t) = 0.399 \times 0.01 \times 100 \sqrt{t}\,(5 - t)$. Elle culmine en $t = 5/3$ ans, soit un tiers de la durée de vie, à environ 1.72 pour un notionnel de 100.

À l'instant $t$, le taux a varié d'une variable normale d'écart-type $\sigma\sqrt{t}$, et la valeur du swap est cette variation multipliée par l'annuité restante, $100(5 - t)$. L'espérance de la partie positive d'une normale centrée vaut $\sigma_X/\sqrt{2\pi} = 0.399\,\sigma_X$.

En annulant la dérivée de $\sqrt{t}\,(5 - t)$, on obtient $(5 - t)/(2\sqrt{t}) = \sqrt{t}$, donc $t = 5/3$. La valeur vaut $0.399 \times 1 \times 1.291 \times 3.333 = 1.72$.

Deux effets se combinent. Le taux a davantage de temps pour bouger, ce qui croît comme $\sqrt{t}$. Et il reste moins de flux, ce qui décroît linéairement.

\iqwrong{Dire que le pic se situe au milieu de la durée de vie. Pour une racine carrée face à une droite, il se situe au tiers.}
\iqfollow{Qu'en serait-il d'un cross-currency swap~?}{Son notionnel est échangé à l'échéance, si bien que l'exposition continue de croître avec $\sqrt{t}$ jusqu'à l'échéance et ne présente pas un tel pic.}
\end{solution}

\begin{solution}{iq:iv2:credit-and-xva:11}
La CVA vaut environ 0.0714 pour un notionnel de 100, soit environ 7 centimes pour 100.

$\mathrm{CVA} = (1 - R)\int_0^5 \mathrm{EE}(t)\, h\,\mathrm{d}t$, et l'intégrale du profil vaut
\[
0.399\left(5 \cdot \tfrac{2}{3}\, 5^{3/2} - \tfrac{2}{5}\, 5^{5/2}\right) = 5.947.
\]
Ainsi $\mathrm{CVA} = 0.6 \times 0.02 \times 5.947 = 0.0714$.

Comme $(1 - R)h$ est le spread de crédit $s = 1.2~\%$, et que l'exposition moyenne vaut $5.947/5 = 1.189$, cela revient à $s \times \overline{\mathrm{EE}} \times T = 0.012 \times 1.189 \times 5$. C'est la règle qu'il faut énoncer à voix haute~: on paie une prime de CDS sur son exposition moyenne.

\iqwrong{Utiliser l'exposition maximale de 1.72 à la place de la moyenne. Cela surestime ici d'environ 45~\%. L'exposition maximale est un indicateur de limite, et non un prix.}
\iqfollow{Que donnerait une actualisation à 3~\% avec pondération par la survie~?}{Elle ramène la valeur à environ 0.064, car l'exposition tardive compte moins.}
\end{solution}

\begin{solution}{iq:iv2:credit-and-xva:12}
Sans netting, il vaut 0.798. Avec netting, il vaut $0.399\sqrt{2 + 2\rho}$~: 0, 0.399, 0.564, 0.691 et 0.798 pour ces cinq corrélations.

Sans netting, chaque transaction contribue pour $0.399$ à elle seule. Avec netting, l'exposition porte sur la somme, qui est normale de variance $2 + 2\rho$. L'espérance de la partie positive vaut 0.399 fois son écart-type.

Le netting n'aggrave jamais l'exposition, car $\max(a + b, 0) \le \max(a, 0) + \max(b, 0)$. À une corrélation de 1, il n'apporte aucun gain. À $-1$, les deux transactions se compensent.

\iqwrong{Dire que des transactions non corrélées ne tirent aucun bénéfice du netting. Elles en tirent beaucoup~: 0.564 contre 0.798, soit environ 29~\% de moins. La corrélation n'a pas besoin d'être négative.}
\iqfollow{Qu'est-ce qui rend le netting valable~?}{Un accord-cadre juridiquement exécutoire conclu avec cette contrepartie. Sans lui, un liquidateur pourrait exiger ce que vous devez tout en vous obligeant à attendre votre tour pour ce qui vous est dû.}
\end{solution}

\begin{solution}{iq:iv2:credit-and-xva:13}
Elle vaut environ 3.99~millions de dollars sans collatéral et environ 0.79~million de dollars avec une marge quotidienne. Elle n'est pas nulle, car le marché bouge pendant les 10 jours qui séparent le dernier appel de marge du close-out.

Sans collatéral, $0.399 \times \$10$ million $\times \sqrt{1}$. Avec marge, seul le mouvement pendant la période de marge en risque compte~: $0.399 \times 10 \times \sqrt{10/252} = 0.795$.

Le rapport vaut donc $\sqrt{10/252} \approx 0.2$. Diviser la période par deux n'économise que 29~\%.

Par ailleurs, la période de marge est souvent plus longue qu'on ne le pense~: contestations, défaut de versement de la marge, puis remplacement de la couverture sur un marché sous tension. Et l'exposition sur cette période tend à se concentrer en période de tension.

\iqwrong{Dire que le collatéral supprime le risque de crédit de contrepartie. Ici il le réduit d'un facteur d'environ cinq, mais pas jusqu'à zéro. La marge initiale est là pour couvrir cet écart.}
\end{solution}

\begin{solution}{iq:iv2:credit-and-xva:14}
Sans pondération, la CVA vaut 0.0714 et la DVA 0.0357, si bien que l'ajustement bilatéral vaut 0.0357 à votre détriment. Avec une pondération first-to-default en $e^{-(h_c + h_o)t}$, elles valent 0.0670 et 0.0335, soit 0.0335.

Un swap à la monnaie a un profil symétrique, si bien que l'exposition négative a la même intégrale, 5.947. Donc $\mathrm{CVA} = 0.6 \times 0.02 \times 5.947$, et $\mathrm{DVA} = 0.6 \times 0.01 \times 5.947$. La pondération ne retient que le premier des deux défauts, puisqu'après celui-ci le contrat a pris fin.

Si votre propre crédit se dégrade, la DVA augmente et vous enregistrez un gain. Une entreprise en difficulté enregistre un gain parce qu'elle va mal.

Les desks s'en méfient pour de bonnes raisons~:
\begin{itemize}
\item On ne peut la monétiser qu'en faisant défaut, ou en vendant une protection sur soi-même, ce qui n'est pas possible.
\item On ne peut la couvrir qu'avec un proxy, comme un panier de pairs, et cela ajoute du risque.
\item Les règles prudentielles la retirent des fonds propres réglementaires, si bien qu'elle ne soutient pas le bilan.
\end{itemize}

La fair value comptable l'inclut encore, car un prix de sortie refléterait votre propre crédit.

\iqwrong{Écarter la DVA comme simplement fictive. Si vous cotez un prix à une contrepartie plus sûre que vous, elle en exigera la prise en compte, si bien que la DVA est réelle dans le prix, même si elle est difficile à monétiser.}
\iqfollow{Et si les deux intensités de défaut sont égales et que le swap est à la monnaie~?}{Alors la CVA et la DVA s'annulent exactement, et vous ne coteriez aucune charge de crédit, avec ou sans pondération.}
\end{solution}

\begin{solution}{iq:iv2:credit-and-xva:15}
Le coût de financement à 100 points de base vaut environ 0.0595 pour un notionnel de 100. Lorsque le spread n'est que du crédit, à 60 points de base, le bénéfice de financement vaut 0.0357, soit exactement votre DVA. Compter les deux revient donc à compter deux fois votre propre crédit.

\textbf{FVA.} Lorsque le swap du client est dans la monnaie pour vous, la couverture est hors de la monnaie, et vous remettez du collatéral dessus. Vous empruntez pour financer ce collatéral, et le client ne remet rien. Vous payez donc le spread de financement sur l'exposition positive~: $0.01 \times 5.947 = 0.0595$.

Lorsque le swap du client est hors de la monnaie, vous recevez du collatéral sur la couverture, ce qui vous évite d'emprunter. Vous gagnez donc le spread sur l'exposition négative, c'est-à-dire le bénéfice de financement.

\textbf{Le chevauchement.} Avec un spread de financement $s_o = h_o(1 - R) = 0.6~\%$, le bénéfice vaut $s_o \times 5.947 = 0.0357$. Et $\mathrm{DVA} = (1 - R)h_o \times 5.947$, le même nombre. C'est une seule chose vue deux fois~: vous pourriez ne jamais rembourser cette dette, ce qui explique à la fois que votre financement coûte plus cher et que votre passif vaut moins.

Un prix cohérent prend donc le coût du côté du financement, et le bénéfice soit du côté de la DVA, soit du côté du financement, mais pas des deux.

Le débat porte sur la question de savoir si la FVA a sa place dans un prix, puisqu'un argument pur d'absence d'arbitrage veut que le financement ne change pas la valeur. Les banques la comptabilisent parce qu'elles la paient. Si le spread de financement est nul, la FVA est nulle, comme il se doit.

\iqwrong{Additionner la CVA, la DVA, le coût de FVA et le bénéfice de FVA sans chercher de chevauchement. Si les spreads sont tous de crédit, la somme compte deux fois votre propre défaut.}
\iqfollow{Le coût de financement dépend-il de la contrepartie~?}{Seulement à travers la survie~: le besoin de financement cesse si le client fait défaut. C'est donc surtout votre propre spread qui compte, ce qui explique que deux banques valorisent différemment la même transaction.}
\end{solution}

\begin{solution}{iq:iv2:credit-and-xva:16}
L'exposition attendue sachant le défaut vaut 1.99 pour $\rho = 0$, 4.25 pour $0.3$, 6.22 pour $0.5$, et 0.62 pour $-0.3$. Ainsi, $\rho = 0.3$ fait plus que doubler la CVA, et $\rho = 0.5$ la triple.

L'exposition est élevée lorsque $Z$ est faible, et avec un $\rho$ positif, le défaut est aussi plus probable lorsque $Z$ est faible. On perd donc le plus lorsque la contrepartie est la plus susceptible de faire défaut~: il s'agit du \index{risque de corrélation défavorable}risque de corrélation défavorable (wrong-way risk).

Sachant le défaut, le facteur est décalé. Le seuil de défaut pour 2~\% vaut $-2.05$, et la moyenne d'une loi normale en dessous de ce seuil vaut $-\varphi(c)/\Phi(c) = -2.42$. Ainsi $Z$ a une moyenne conditionnelle de $-2.42\rho$, et $E[\max(-5Z, 0) \mid \text{défaut}]$ est une intégrale~:
\[
\frac{1}{0.02}\int_{-\infty}^{0} (-5z)\,\varphi(z)\,\Phi\!\left(\frac{c - \rho z}{\sqrt{1 - \rho^2}}\right)\mathrm{d}z.
\]
Pour $\rho = 0$, elle vaut $5\varphi(0) = 1.99$, soit le chiffre de référence.

\textbf{Exemples de wrong-way risk.} Acheter une protection sur une banque auprès d'une autre banque du même système financier. Détenir une créance libellée en dollars sur un emprunteur dont la devise est faible, celle-ci baissant au moment où il fait défaut. Prêter contre les actions émises par la contrepartie elle-même.

\textbf{Exemple de right-way risk.} Une compagnie aérienne qui a acheté du carburant à terme auprès de vous. Vous êtes exposé lorsque le carburant baisse, et cette baisse avantage la compagnie. Il faut toujours vérifier le sens~: un producteur de pétrole qui a acheté du pétrole à terme auprès de vous est en wrong-way risk, car la baisse du pétrole fait croître votre exposition et fragilise le producteur.

\iqwrong{Multiplier l'exposition moyenne par la probabilité de défaut et s'arrêter là. Cela suppose l'indépendance, et l'erreur peut atteindre un facteur trois.}
\end{solution}

\begin{solution}{iq:iv2:credit-and-xva:17}
Le tableau donne la perte attendue, exprimée en part de chaque tranche.

\begin{center}
\begin{tabular}{lcccc}
\toprule
$\rho$ & 0--3~\% & 3--7~\% & 7--15~\% & 15--100~\% \\
\midrule
0.05 & 39.6~\% & 0.33~\% & 0.00~\% & 0.00~\% \\
0.2 & 33.5~\% & 4.06~\% & 0.40~\% & 0.00~\% \\
0.5 & 21.3~\% & 6.70~\% & 2.42~\% & 0.12~\% \\
\bottomrule
\end{tabular}
\end{center}

La tranche equity perd de la valeur quand la corrélation baisse et en gagne quand elle monte. La tranche senior évolue en sens inverse. Ce qui reste fixe, c'est la perte attendue du pool, $0.02 \times 0.6 = 1.2~\%$, puisqu'une somme d'espérances ignore la corrélation.

Conditionnellement au facteur de marché $M$, la perte d'un grand pool est déterministe~:
\[
L(M) = (1 - R)\,\Phi\!\left(\frac{\Phi^{-1}(p) - \sqrt{\rho}\,M}{\sqrt{1 - \rho}}\right).
\]
La perte de la tranche vaut $E[\min(\max(L - a, 0), b - a)]/(b - a)$, obtenue par quadrature sur $M$.

Il y a deux limites. Quand $\rho \to 0$, le pool perd exactement 1.2~\%, donc la première tranche a perdu $1.2/3 = 40~\%$, et le reste rien. Quand $\rho \to 1$, tous les défauts surviennent ensemble avec une probabilité de 2~\%, et la perte vaut 60~\%, si bien que chaque tranche jusqu'à 60~\% a perdu 2~\%.

La tranche equity est donc longue de corrélation, et la senior est courte de corrélation. La tranche mezzanine évolue dans un sens puis dans l'autre~: elle monte ici, mais doit redescendre vers 2~\% quand $\rho \to 1$.

\iqwrong{Penser qu'il existe une seule corrélation pour toutes les tranches du marché. Une corrélation de modèle cotée diffère d'une tranche à l'autre. C'est le smile de corrélation, et il montre que la copule gaussienne à un facteur ne s'ajuste pas.}
\iqfollow{Où ce modèle échoue-t-il en période de crise~?}{Il n'a pas de dépendance de queue, et la corrélation de chaque paire est un nombre unique. L'usage que Li a fait de la copule gaussienne a rendu la corrélation de défaut facile à coter, mais il ne dit pas que les défauts se regroupent davantage en période difficile.}
\end{solution}

\begin{solution}{iq:iv2:credit-and-xva:18}
La valeur des actions vaut 24.15, celle de la dette 75.85, la probabilité de défaut risque-neutre est de 18.7~\%, et le spread de crédit vaut environ 232 points de base. À trois mois, le spread est d'environ 71 points de base, et à cinq ans d'environ 210.

Les actions sont un call sur les actifs. On a donc $d_1 = 1.1376$, $d_2 = 0.8876$, et $E = 100\,\Phi(d_1) - 80 e^{-0.03}\Phi(d_2) = 24.15$. La dette est ce qui reste~: $100 - 24.15 = 75.85$. Une seconde voie consiste à voir la dette comme une obligation sans risque moins un put~: $77.64 - 1.78 = 75.85$.

Il y a défaut quand les actifs finissent sous 80, donc la probabilité risque-neutre vaut $\Phi(-d_2) = 18.7~\%$.

Le spread vaut $-\ln(75.85/80) - 0.03 = 2.32~\%$.

\textbf{Structure par terme.} À un mois, le spread vaut environ 2 points de base, et il tend vers zéro quand la maturité tend vers zéro. Il culmine vers deux ans à environ 248, puis baisse~: 210 à cinq ans, 162 à dix ans.

L'extrémité courte est le point faible. Une entreprise ne peut atteindre le défaut que par une diffusion continue, et sur un horizon court elle n'y parvient pas. Le modèle donne donc presque rien pour le crédit à court terme, alors que les spreads courts réels ne sont pas proches de zéro. Les modèles ajoutent des sauts, ou une incertitude sur le niveau des actifs ou sur la barrière, pour corriger cela.

Par ailleurs, la valeur des actifs et leur volatilité ne sont pas observées. Les praticiens les déduisent des actions et de la volatilité des actions, en utilisant $\sigma_E E = \Phi(d_1)\sigma_A V$.

\iqwrong{Lire 18.7~\% comme une probabilité réelle. Le calcul utilise la dérive risque-neutre. Avec une dérive réelle des actifs supérieure au taux, la vraie probabilité est plus faible.}
\iqfollow{Qu'est-ce que la distance au défaut~?}{C'est $(\ln(V/D) + (\mu - \sigma^2/2)T)/(\sigma\sqrt{T})$, le nombre d'écarts-types qui séparent les actifs de la dette. C'est $d_2$ avec la dérive réelle à la place du taux.}
\end{solution}

\begin{solution}{iq:iv2:credit-and-xva:19}
Le triangle du crédit donne 1.67~\%, 2.92~\% et 4.58~\%. Le bootstrap donne 1.67~\%, 2.95~\% et 4.81~\%.

\textbf{Triangle.} Chaque spread donne une intensité de défaut moyenne~: $0.01/0.6$, $0.015/0.6$, $0.02/0.6$. Les morceaux forward se déduisent de l'intensité totale~: $(3 \times 2.5~\% - 1.67~\%)/2 = 2.92~\%$ et $(5 \times 3.33~\% - 3 \times 2.5~\%)/2 = 4.58~\%$.

\textbf{Bootstrap.} On résout dans l'ordre. On prend d'abord l'intensité de la première période, de sorte que le contrat à un an vaille zéro. Puis, en la gardant, on prend la deuxième pour que le contrat à trois ans vaille zéro, puis la troisième. Jambe de prime~: $\sum_i \tfrac14 D_i S_i$ plus une demi-période de couru en cas de défaut. Jambe de protection~: $(1 - R)\sum_i D_i (S_{i-1} - S_i)$.

\textbf{Pourquoi plus élevé.} Un spread de CDS est une moyenne pondérée des spreads forward, et les poids sont les facteurs d'actualisation risqués $D_iS_i$, qui ne sont pas égaux. Les premières périodes pèsent davantage, car elles sont moins actualisées et plus probablement atteintes. Les poids temporels égaux du triangle sous-estiment donc à quel point le dernier morceau doit être élevé pour relever la moyenne à 200.

\textbf{La courbe inversée.} À 2\,000 et 800, la première année a une intensité d'environ 33~\%. L'intensité totale sur cinq ans par le triangle vaut $5 \times 13.3~\% = 66.7~\%$, donc les quatre années suivantes en moyenne environ 8.3~\%. Le marché dit ceci~: ce nom est en danger immédiat, et s'il s'en sort, il sera beaucoup plus sûr. Le contrat long d'un nom en difficulté est proche d'une créance sur le même événement à court terme.

\iqwrong{Ajuster une intensité plate unique à chaque maturité et appeler cela la courbe. Des intensités plates différentes selon les maturités ne sont pas cohérentes avec une seule courbe de survie, et elles ne permettent pas d'évaluer un contrat à terme à départ différé.}
\iqfollow{Qu'est-ce qui peut mal tourner dans un bootstrap~?}{Une courbe fortement inversée avec un taux de recouvrement élevé peut forcer une intensité forward négative, ce qui est un arbitrage dans les données d'entrée ou le signe que l'hypothèse de recouvrement est trop élevée.}
\end{solution}

\begin{solution}{iq:iv2:credit-and-xva:20}
Le CVA vaut environ 110\,800~\$. Sa sensibilité au spread vaut environ 515~\$ par point de base, donc il faut acheter environ 1.1~million de dollars de protection. Après le mouvement de taux, le CVA vaut 182\,500~\$. Après le mouvement de spread, il vaut 160\,600~\$. Après les deux, il vaut 265\,400~\$. Les deux effets ensemble ajoutent donc 33\,100~\$ au-delà de la somme.

\textbf{Valeur.} Avec $h = 0.02/0.6 = 3.33~\%$ et $\mathrm{EE}(t) = 0.399\sqrt{t}\,(5 - t)$ pour 100,
\[
\mathrm{CVA} = 0.6 \int_0^5 h e^{-ht}\,\mathrm{EE}(t)\,\mathrm{d}t = 0.1108 \text{ per } 100.
\]
Sans la pondération par la survie, il vaudrait $0.02 \times 5.947 = 0.119$.

\textbf{Couverture.} On réévalue avec le spread à 201~: le CVA augmente de 515~\$. Un CDS de 1~million de dollars à cinq ans gagne son annuité risquée multipliée par un point de base~: $(1 - e^{-5h})/h \times 10^{-4} \times \$1$ million $= 4.61 \times \$100 = \$461$. Il faut donc $515/461 = \$1.12$ million de protection, soit à peu près 1~\% du notionnel du swap. Le montant est faible parce que l'exposition est petite face au notionnel.

\textbf{Mouvement de taux.} Si le taux a bougé de sorte que le swap soit favorable à la position de 50 points de base, l'exposition a une moyenne $m = 0.005$ et
\[
\mathrm{EE}(t) = 100(5 - t)\,\sigma_t\,\bigl(\varphi(m/\sigma_t) + (m/\sigma_t)\Phi(m/\sigma_t)\bigr), \quad \sigma_t = \sigma\sqrt{t}.
\]
Cela donne 182\,500~\$. Un mouvement de 50 points de base défavorable au swap donnerait 64\,100~\$, donc la sensibilité n'est pas symétrique, et le delta à la monnaie vaut environ 1\,180~\$ par point de base.

\textbf{Les deux.} Additionner les mouvements séparés donne $110\,800 + 71\,700 + 49\,700 = \$232\,200$, alors que la vraie valeur est de 265\,400~\$. Les 33\,100~\$ supplémentaires sont le \index{gamma croisé}gamma croisé entre taux et crédit. Quand on est dans la monnaie, le CS01 vaut 855~\$, contre 515~\$. La couverture en spread doit donc faire 1.86~million de dollars, et elle était insuffisante de 0.74~million de dollars pour exactement le mouvement qui fait mal.

Il faut donc rééquilibrer la couverture de crédit quand les taux bougent, et la couverture de taux quand les spreads bougent. Et un desk qui couvre chaque risque séparément paie quand même le terme croisé. Avec un risque de corrélation défavorable, le terme croisé grandit, car spread et exposition évoluent alors ensemble.

\textbf{Pas de CDS négocié.} On couvre avec un proxy, comme un indice sectoriel et de notation ou un nom apparenté. On retire le risque de spread systématique, et on garde le risque idiosyncratique. C'est un résidu réel, et on détient des réserves ou du capital pour le couvrir. On indique clairement que le saut au défaut n'est pas couvert.

\begin{dated}{2026-11}{Pourquoi existe la charge de CVA}\label{dat:iv2:credit-and-xva:cva}
Le communiqué de presse du Comité de Bâle sur le traitement en capital du risque de contrepartie bilatéral indique que, pendant la crise, environ deux tiers des pertes attribuées au risque de crédit de contrepartie provenaient de pertes de CVA, et seulement un tiers environ de défauts effectifs. C'est pourquoi les règles de fonds propres imposent une charge sur la variation de valorisation au prix de marché du CVA, et pas seulement sur le défaut.
\end{dated}

\iqwrong{Couvrir le CVA seulement pour le spread, comme un seul nombre. Le CVA est une position de type option sur la valorisation au prix de marché du portefeuille, multipliée par une courbe de crédit. Il présente un delta de taux, un vega et un gamma croisé, et c'est de là que viennent les pertes.}
\iqfollow{Une hausse de 20~\% de la volatilité des taux a-t-elle de l'importance~?}{À la monnaie, oui, presque en proportion~: le CVA passe à 133\,000~\$, car l'exposition est linéaire en volatilité. Un desk de CVA est donc court de vega sur les taux de son livre~: il perd quand la volatilité monte.}
\end{solution}
